\section{Setting and theory}\label{sec:setting}

\subsection{Captive spin-offs in Korea}\label{sec:captive}

Korean company law allows a listed company to divide itself in two ways. In an equity spin-off (\emph{injŏk bunhal}) the new company's shares are distributed to the parent's shareholders in proportion to their holdings, as in a US spin-off; in a captive spin-off (\emph{mulchŏk bunhal}) the parent itself receives all of the new company's shares, so that the parent's shareholders now own the carved-out business only through the parent \cite{kang2025}. The second form has dominated: of the \KLKrxTotal{} division announcements the Korea Exchange counted from 1998, when the Commercial Act first allowed divisions, to 2022, \KLKrxCaptive{} were captive \cite{kang2025}; \citeA{kang2025a} analyse the choice between the two forms and their announcement effects. The form has no immediate effect on the consolidated statements and none on the parent's ownership structure. Its consequences arrive later, when the parent decides what to do with the subsidiary it now wholly owns. It may keep it, sell part or all of it, or list it, and each of these decisions is taken by the parent's board, which the controlling owner appoints, without a vote of the parent's shareholders.

The disposal is what the public controversy has been about. When LG Chem announced in September 2020 that it would carve out its battery business into a wholly owned subsidiary, its filing gave operational efficiency, specialisation, the subsidiary's competitiveness and its long-term growth as the purposes. The shareholders' meeting approved the division on 30 October 2020 over the opposition of some large outside investors; the subsidiary was established on 1 December 2020 as LG Energy Solution and listed on 27 January 2022, by which date LG Chem's stock had lost more than \KLLGDrop{} percent of its value since the announcement \cite{kang2025}. In our data LG Chem's filing carries a disposal signal of \LGScore{} on a scale from zero to one; the stock's three-day announcement return was \LGAnnCar{} percent and its market-adjusted return over the following year \LGBhar{} percent. By 2022 the protection of general shareholders in captive spin-offs was an item on the new administration's list of national policy tasks, and the Financial Services Commission spent several months consulting experts and stakeholders before it acted \cite{fsc2022a}. The objection practitioners raised, in the press and at shareholders' meetings, has two parts that matter for what follows. The first is that the loss to the parent's minority shareholders occurs at the disposal, not at the division. The second is that the controlling owner knows, at the division, more about the disposal than the filing lets the public know.

The regulator answered in 2022. On 5 September 2022 the Financial Services Commission announced three measures for captive spin-offs \cite{fsc2022a}. First, the material-event report that announces the division must state the division's specific purpose (restructuring, sale or listing, for example), its expected effects and the measures the company will take to protect its shareholders, and must disclose and keep current any plan to list the new company. Second, shareholders who vote against a captive spin-off receive an appraisal right to sell their shares to the company at a price based on the market price before the division was pursued. Third, when a carved-out subsidiary applies for listing within five years of the division, the Korea Exchange reviews how the parent protected its general shareholders and may refuse the listing. The listing review applied from 28 September 2022 and the revised disclosure form from 18 October 2022. The appraisal right required an amendment to the enforcement decree of the Capital Markets Act, which the Cabinet approved on 20 December 2022 for promulgation and immediate effect before the end of the year \cite{fsc2022b}. Of the captive spin-offs in our data announced from 2017 to the September 2022 package, none states shareholder-protection measures; the first that does is dated 23 September 2022, between the package and the revised form, and \ProtectPost{} percent of the filings received from October 2022 onward do.

Two features of this setting make it useful for the question we ask. The disposal is a discrete, observable event that the parent must report in its annual report, so the controlling owner's revealed intention can be measured for every division, years later, from public filings. And the filing that announces the division is a long, structured document (the parent must state the purpose of the division, the business and capital of the new company, the division ratio, whether the new company will apply for listing, and from October 2022 the protection measures), so what a reader could have inferred at the announcement can be measured as well. Together they let us ask whether the controlling owner's intention was legible when minority shareholders still held the asset, and what the market did with that information.

\subsection{Controlling owners and the disposal decision}\label{sec:theory}

We read the captive spin-off through the principal--principal frame in which the governance problem of concentrated-ownership firms is the conflict between controlling and minority owners rather than between owners and managers \cite{young2008}. The frame has a specific prediction for Korea. Controlling families hold the parent through pyramids and cross-holdings, so their cash-flow stake in a listed parent is often far below their control stake \cite{laporta1999, claessens2000}. A transfer of value out of the parent---tunneling, in the term of \citeA{johnson2000}---then costs the family only its cash-flow share of the loss, while the gain, if it accrues elsewhere in the group or to the family directly, is theirs in full. Firm value falls as the wedge between control and cash-flow rights widens \cite{claessens2002, masulis2011}. Korean business groups have been the canonical setting for documenting such transfers: mergers that favour the controlling family's firms \cite{bae2002}, private securities offerings priced in the family's favour \cite{baek2006}, propping within groups when a member is in trouble \cite{bae2008}, and lower profitability and larger crisis losses where the controller's cash-flow stake was small \cite{joh2003, mitton2002, black2006}. The structure of the groups themselves is shaped by the controller's position in the pyramid \cite{almeida2011}. Outside Korea the same arithmetic appears in intercorporate loans in China \cite{jiang2010} and in the premia paid for control blocks across countries \cite{dyck2004}. The most recent Korean evidence in this journal is of the same kind: when a controlling family incurs an inheritance or gift tax, its firms raise dividends, and more so where the potential for expropriation is high \cite{na2025}. The reviews of this literature make the same point in general form: where the controller's cash-flow stake is small, the wedge between control and ownership is the governance variable to watch \cite{young2008, castaner2022, aguilera2025}.

A captive spin-off is a transfer waiting to happen. The division itself moves nothing; it changes the set of decisions the controller can take without the minority. After the division the controller can sell the subsidiary to a third party, sell it to another group company, bring in an investor, or list it, and in each case the minority shareholders of the parent receive their share of the proceeds only insofar as the parent's value reflects them. \citeA{kang2025} document, for \KLN{} captive spin-offs announced between 2001 and 2021, that the parent's stock earns a small positive return at the announcement and a large negative abnormal return over the following year, that the loss is concentrated where the controlling owner's stake is small, and that it is concentrated in the divisions whose subsidiary was eventually disposed of. They also write, without testing it, that investors are unlikely to discern the controller's intention at the announcement. That sentence is our starting point.

Our first hypothesis concerns the filing. Its premise is that a division filing is written by the parent, under the controller's direction, for a purpose; a controller who intends to list or sell the subsidiary has reasons to say so, or to say things that imply it (a stated purpose of attracting investment, a reference to a listing application, a description of the new company as a vehicle for outside capital), because the filing is also read by the investors and buyers the controller will need. The mechanism is therefore not a leak but ordinary disclosure: intentions shape the stated purpose, the description of the new company and the choice of what to leave out. The boundary of the hypothesis is that the filing can only be informative about intentions that exist at the announcement; disposals decided years later for reasons that arose later are noise against which any reader, human or machine, will fail. The rival prediction is that the filing is boilerplate, written to a form, so that a careful reading predicts the disposal no better than the parent's financial position does. We therefore expect (H1) that a reading of the filing text predicts whether the subsidiary is later disposed of, and does so better than the parent's and the new company's financial characteristics at the time of the division.

The second hypothesis concerns the controller's stake. The premise is the principal--principal arithmetic: the controller's gain from a disposal (listing proceeds that fund the group, a sale to a related company, a stake sold to an investor on terms the controller sets) does not depend on the controller's cash-flow stake in the parent, while the controller's share of the parent's loss does. The mechanism is that a smaller stake lowers the controller's private cost of a disposal and so raises the probability that the division is undertaken with a disposal in view. The boundary is that the stake is one determinant among several; a controller with a large stake may still dispose because the subsidiary needs capital the parent cannot supply, and a controller with a small stake may retain because the subsidiary is the group's core. The rival prediction comes from the efficiency view of divestitures, in which the division is a refocusing or financing device \cite{krishnaswami1999, slovin1995, powers2003, vijh2002}: the controller's stake is unrelated to the disposal because the disposal is decided on the subsidiary's merits. \citeA{allen1998}, who find that US carve-outs by parents that keep the proceeds are received worse than those that pay them out, supply the version of this view in which the market reads the parent's intentions at the announcement. We expect (H2) that the smaller the controlling block before the announcement, the stronger the disposal signal in the filing and the more likely the subsidiary is disposed of, particularly soon after the division.

The third hypothesis concerns the market, and here we state the hypothesis together with the rival that \citeA{kang2025} assert. If minority investors read the filing as a careful reader would, the premise of efficient pricing implies two things. A filing which signals a disposal should be met with a weaker, or more negative, announcement return than one that signals retention, as investors in family firms price a death in the controlling family as a dilution of control, the more so where shareholder rights are weak \cite{tanyerigunsur2023}. And nothing systematic should follow: the loss from the disposal would be in the price at the announcement, and the subsequent return would be unrelated to the signal. The mechanism is the ordinary incorporation of public information. The boundary is that the test requires the signal to be measurable at the announcement, which H1 supplies, and that the one-year window captures the disposals that follow quickly. The rival prediction is the one practitioners and \citeA{kang2025} make. Investors do not discern the intention; the announcement return is unrelated to the signal (or, if the division is read as a refocusing, positively related to a stated growth purpose); and the loss arrives afterwards as the disposal materialises, so that the subsequent return falls with the signal and with the realised disposal. We state (H3) as the efficient-reading hypothesis: a more negative announcement return and no subsequent drift where the filing signals disposal. The opposite pattern---a positive or null announcement reaction followed by a negative drift concentrated in the parents that dispose early and in the parents with small controlling blocks---is evidence for the rival.

The fourth hypothesis concerns the reform. The premise is that mandated disclosure acts on what is written, and on the price reaction to what is written, before it acts on what is done. The mechanism by which it could act on what is done is the appraisal right and the listing review, which raise the controller's cost of a disposal-oriented division; the mechanism by which it acts on what is written is the form itself. The boundary is time. The appraisal right took effect in late December 2022, so the first post-reform quarter had the new form without the new right; and disposals take a year or more to be observed, so the post-reform disposal record is short. The rival prediction, in which the reform deters disposal-oriented divisions, is that the disposal signal in post-reform filings is weaker and realised early disposal rarer. We expect (H4) that after the reform filings state protection measures and listing intentions where they did not before, that the announcement reaction changes, and that the disposal signal and the rate of early disposal do not, over the horizon we can observe.

These are predictions about associations in public records. We do not have an instrument for the controlling block or for the division decision, and we do not claim that the block causes the disposal or the return. We claim that the pattern of associations is the one the principal--principal frame predicts and not the one the efficiency view predicts, and we say where the data cannot tell the two apart.
